\documentclass{article}

\PassOptionsToPackage{numbers, sort, compress}{natbib}
\usepackage[main,final]{neurips_2025}

\usepackage{fontspec}
\setmainfont{Times New Roman}
\usepackage[russian]{babel}
\usepackage{hyperref}
\usepackage{url}
\usepackage{booktabs}
\usepackage{nicefrac}
\usepackage{microtype}
\usepackage{xcolor}

%%%

\usepackage{subcaption}
\usepackage{graphicx}
\usepackage{multirow}
\usepackage{amsmath,amssymb,amsfonts}
\usepackage{amsthm}
\usepackage{mathrsfs}
\usepackage{xcolor}
\usepackage{textcomp}
\usepackage{manyfoot}
\usepackage{algorithm}
\usepackage{algorithmicx}
\usepackage{algpseudocode}
\usepackage{listings}
\hypersetup{colorlinks=true,linkcolor=black,citecolor=black,urlcolor=black}

% Русский заголовок аннотации; служебный колонтитул шаблона отключён.
\renewenvironment{abstract}{%
  \begin{center}\large\bfseries Аннотация\end{center}\begin{quote}%
}{\end{quote}}
\makeatletter
\renewcommand{\@notice}{}
\makeatother

\newtheorem{theorem}{Теорема} % сквозная нумерация
%%\newtheorem{theorem}{Теорема}[section] % нумерация внутри раздела
%% Необязательный аргумент [theorem] задаёт общую нумерацию теорем и утверждений.
\newtheorem{proposition}[theorem]{Утверждение}%
\newtheorem{lemma}{Лемма}%
%%\newtheorem{proposition}{Утверждение} % для отдельной нумерации утверждений

\newtheorem{example}{Пример}
\newtheorem{remark}{Замечание}

\newtheorem{definition}{Определение}
\newtheorem{assumption}{Предположение}

%%%


\title{Влияние моделей источников (быстрых БЯМ) на качество декодирования текста при помощи ЭЭГ-спеллера}


% Команда \author поддерживает любое число авторов. \And и \AND разделяют
% имена: \And оставляет выбор переноса LaTeX, \AND задаёт принудительный перенос.


\author{
  Барабанщиков М.~А.\\
  автор\\
  \And
  Майсурадзе А.~И.\\
  научный руководитель и консультант\\
}


\begin{document}


\maketitle

\begin{abstract}
Исследуется влияние моделей источника на декодирование текста в ЭЭГ-спеллере на основе вызванного потенциала P300. Языковая модель задаёт априорное распределение следующего символа по уже введённому тексту, а ЭЭГ-классификатор оценивает свидетельства о выборе при предъявлении стимулов. Рассматриваются статистическая модель и локальная быстрая большая языковая модель, а также две стратегии предъявления: статическая с фиксированным набором символов и динамическая, изменяющая набор кандидатов с сохранением доступа ко всем символам. Задача формулируется как минимизация ошибки декодирования и времени ввода при ограничении на задержку системы. Предлагается сравнить конфигурации при общем классификаторе ЭЭГ и единых разбиениях данных; качество предполагается оценивать как в офлайн-режиме, так и при онлайн-вводе текста.
\end{abstract}

\textbf{Ключевые слова:} ЭЭГ-спеллер, интерфейс мозг-компьютер, P300, языковая модель, декодирование текста.

\section{Введение}\label{sec:intro}

Интерфейсы мозг-компьютер позволяют вводить текст без обычного движения рук или речи, что особенно важно для людей с тяжёлыми двигательными нарушениями. В ЭЭГ-спеллере пользователь выбирает символ, сосредоточивая внимание на зрительном стимуле; вызванный потенциал P300 в электроэнцефалограмме помогает определить намерение. Однако для каждой буквы приходится предъявлять стимулы и накапливать достаточно свидетельств, поэтому ошибки распознавания и длительное время выбора ограничивают практическую скорость общения. Цель этой работы — выяснить, насколько контекстно-зависимая модель источника может улучшить декодирование текста при приемлемой вычислительной задержке.

Модель источника задаёт априорную вероятность следующего символа по уже набранному тексту; классификатор ЭЭГ оценивает свидетельство о выборе на основе сигнала. Их сочетание позволяет ранжировать кандидатов с учётом обеих частей информации. В PolyMorph контекст использовался для изменения матрицы выбора и предсказания слов \citep{casagrande2015polymorph}; позднее GPT-2 применяли для оптимизации стимулов и подсказок в P300-спеллере \citep{parthasarathy2024gpt2}. ChatBCI показал онлайн-ввод с подсказками большой языковой модели через удалённый сервис \citep{hong2026chatbci}, а MindChat исследовал похожие подсказки для SSVEP-спеллера \citep{wang2025mindchat}.

Эти результаты подтверждают пользу языкового контекста, но сопоставлять их количественно трудно: различаются парадигмы стимуляции, состав интерфейса, способ оценки и доступность модели. В частности, выигрыш от подсказок слов не равен выигрышу от более точного априорного распределения по символам, а моделирование ЭЭГ не заменяет онлайн-испытание. Отдельная линия работ переводит ЭЭГ при чтении текста непосредственно в слова: DeWave \citep{duan2023dewave}, метод Wang и Ji \citep{wang2022open} и модель с контрастивным предобучением \citep{wang2024contrastive} решают иную задачу, поэтому их метрики не служат прямой базой для P300-спеллера. Критический анализ \citep{jo2024working} дополнительно показывает необходимость оценки без утечки правильного текста и с контролем вклада нейросигнала.

Предлагается единая конфигурируемая система, в которой один и тот же декодер ЭЭГ можно сочетать со статистической моделью или быстрой БЯМ. Рассматриваются две стратегии: статическая сохраняет фиксированный набор предъявляемых символов и меняет только априорные вероятности; динамическая также адаптирует набор кандидатов на экране. Основной исследовательский вопрос состоит в том, даёт ли динамическая стратегия устойчивый выигрыш по сравнению со статической при сопоставимых условиях, одном классификаторе и одинаковом ограничении на задержку. Система позволяет отделить вклад языковой модели от вклада интерфейса и проверить обе стратегии в одном протоколе; возможная цена динамической стратегии — дополнительное время вычисления и изменение стимуляции.

Для проверки предусмотрены офлайн-сравнение на одинаковых разбиениях ЭЭГ и онлайн-ввод текста участниками. Базовыми конфигурациями служат декодирование без языковой модели, со статистической моделью и с быстрой БЯМ; статическое и динамическое предъявление сравниваются при одинаковом классификаторе ЭЭГ. Критерии оценки включают точность классификации P300 и выбора символа, скорость ввода, частоту исправления ошибок, качество предсказаний и полную задержку отклика. Эффект динамической стратегии необходимо проверять онлайн, поскольку изменение стимулов меняет собираемый ЭЭГ-сигнал и не всегда воспроизводится по записям, полученным при фиксированной раскладке.

\textbf{Предлагаемый вклад.}
\begin{itemize}
    \item Единый протокол сравнения статистических и быстрых нейросетевых моделей источника в P300-спеллере.
    \item Сопоставление статического и динамического предъявления символов при общем декодере ЭЭГ.
    \item Протокол офлайн- и онлайн-оценки точности, скорости и задержки с разделением обучающих и контрольных данных.
\end{itemize}

\textbf{Структура статьи.} Далее рассмотрены связанные работы и приведена формальная постановка задачи.

\section{Обзор литературы}\label{sec:rw}

\textbf{Контекст и адаптация в P300-спеллерах.} Ранний PolyMorph \citep{casagrande2015polymorph} удаляет малополезные символы из матрицы и предсказывает слова по контексту предложения. Это показывает, что языковая информация может менять не только решение декодера, но и сами стимулы. Parthasarathy и соавторы \citep{parthasarathy2024gpt2} совместили GPT-2 с оптимизацией предъявления и предсказанием слов; сообщаемый выигрыш получен преимущественно в моделировании на выборках ЭЭГ. Оба подхода важны как базовые, но не дают прямого сравнения современных быстрых БЯМ и статистических априорных моделей при одном онлайн-протоколе.

\textbf{Подсказки большой языковой модели.} В ChatBCI \citep{hong2026chatbci} слова и продолжения фразы становятся дополнительными клавишами P300-интерфейса. Онлайн-эксперимент с семью участниками подтверждает, что подсказки сокращают число выборов, однако обращения к удалённому GPT-3.5 не отвечают на вопрос о задержке локальной быстрой модели. MindChat \citep{wang2025mindchat} расширяет идею контекстных подсказок до слов и предложений, но использует SSVEP, поэтому его показатели скорости нельзя напрямую сопоставить с P300. Для настоящей работы нужны две независимые переменные: качество распределения вероятностей символов и способ предъявления кандидатов.

\textbf{Декодирование текста непосредственно из ЭЭГ.} Wang и Ji \citep{wang2022open}, DeWave \citep{duan2023dewave} и Wang с соавторами \citep{wang2024contrastive} соединяют нейросигнал и языковые представления в задаче перевода ЭЭГ, записанной при чтении, в текст. Последние две работы опубликованы на NeurIPS и ACL соответственно и показывают развитие методов сопряжения ЭЭГ с языковыми моделями. Однако они не используют выбор символов по P300, а их BLEU и ROUGE описывают восстановление текста, а не скорость ввода. Более того, для ранней модели Wang и Ji была обнаружена ошибка оценки; анализ Jo и соавторов \citep{jo2024working} требует тестировать модель без подстановки правильного продолжения и сравнивать её с контролем на шумовом сигнале.

\textbf{Вывод.} Ближайшие базовые подходы — P300 без языковой модели, P300 со статистическим априорным распределением и P300 с предсказанием текста посредством БЯМ. Пробел состоит в сопоставлении этих подходов на общих данных и в проверке, окупает ли динамическое изменение набора стимулов дополнительные вычисления. Именно эту проверку должен обеспечить предлагаемый эксперимент.

\section{Постановка задачи}\label{sec:problem}

Требуется выбрать намеренный пользователем символ по последовательности предъявленных стимулов, фрагментам ЭЭГ и ранее введённому тексту. Исследовательский вопрос состоит в том, как выбор модели источника и способа предъявления символов влияют на ошибку декодирования и время ввода при одинаковом классификаторе ЭЭГ и заданном ограничении на задержку.

\textbf{Данные и обозначения.} Пусть $\mathcal A$~--- конечный алфавит доступных символов. Выборка $\mathcal D_{\pi}=\{(h_i,y_i,X_i)\}_{i=1}^{n_{\pi}}$, собранная при политике предъявления $\pi$, состоит из контекста $h_i\in\mathcal A^*$ перед очередным выбором, намеренного символа $y_i\in\mathcal A$ и записи $X_i=\{(x_{ij},S_{ij})\}_{j=1}^{m_i}$. Здесь $x_{ij}$~--- фрагмент ЭЭГ после $j$-го предъявления, а $S_{ij}$~--- предъявленная группа клавиш. Целевая метка предъявления равна $z_{ij}=\mathbf 1\{g_{C_i}(y_i)\in S_{ij}\}$. Отображение $g_{C_i}$ оставляет символы текущего набора $C_i\subseteq\mathcal A$ без изменения и переводит остальные в клавишу «Другие». При её выборе открывается полный алфавит; поэтому каждый $y_i$ остаётся доступным и при динамической раскладке.

\textbf{Класс моделей и правило решения.} Модель источника $p_{\phi}(c\mid h)$ задаёт положительную априорную вероятность $c\in\mathcal A$; рассматриваются равномерное распределение без языковой модели, статистическая модель и быстрая БЯМ. Классификатор ЭЭГ с параметрами $\theta$ оценивает калиброванное логарифмическое отношение правдоподобий $s_{\theta}(x)=\log[p_{\theta}(x\mid z=1)/p_{\theta}(x\mid z=0)]$. Политика $\pi$ задаёт набор клавиш и группы стимулов: статическая использует $C_i=\mathcal A$, динамическая выбирает $C_i=\pi(h_i)$ и применяет описанный возврат к полному алфавиту. Для $a\in C_i\cup\{\text{Другие}\}$ априорная масса $p_{\phi}^{(C_i)}(a\mid h_i)$ равна $p_{\phi}(a\mid h_i)$ для символа и сумме вероятностей исключённых символов для клавиши «Другие». После предъявлений оценка выбора имеет вид
\begin{equation}\label{eq:posterior}
q_{\theta,\phi,\pi}(a\mid X_i,h_i)
=\frac{p_{\phi}^{(C_i)}(a\mid h_i)\exp\!\left(\sum_{j=1}^{m_i}\mathbf 1\{a\in S_{ij}\}s_{\theta}(x_{ij})\right)}
{\sum_{b\in C_i\cup\{\text{Другие}\}}p_{\phi}^{(C_i)}(b\mid h_i)\exp\!\left(\sum_{j=1}^{m_i}\mathbf 1\{b\in S_{ij}\}s_{\theta}(x_{ij})\right)}.
\end{equation}
Окончательная вероятность $P_{\theta,\phi,\pi}(y_i\mid X_i,h_i)$ учитывает также второй этап выбора, если была нажата клавиша «Другие»; итоговый символ определяется максимумом этой вероятности. Формула~\eqref{eq:posterior} предполагает условную независимость свидетельств предъявлений при заданном намеренном символе; калибровка $s_{\theta}$ и устойчивость к нарушению этого предположения подлежат проверке.

\textbf{Критерий и ограничения.} Пусть $T_i(\theta,\phi,\pi)$~--- полное время выбора символа, включая повторные предъявления, переход через клавишу «Другие» и вычисление модели. При $\lambda\geqslant0$ внутренний критерий объединяет логарифмическую ошибку и время ввода:
\begin{equation}\label{eq:objective}
\begin{aligned}
Q_{\mathcal D_{\pi}}(\theta,\phi,\pi)
&=\frac1{n_{\pi}}\sum_{i=1}^{n_{\pi}}\left[-\log P_{\theta,\phi,\pi}(y_i\mid X_i,h_i)
  +\lambda T_i(\theta,\phi,\pi)\right],\\
(\theta^*,\phi^*,\pi^*)
&\in\underset{(\theta,\phi,\pi)\in\mathcal H,\;L_{\mathrm{sys}}\leqslant L_{\max}}{\arg\min}
Q_{\mathcal D_{\pi}}(\theta,\phi,\pi).
\end{aligned}
\end{equation}
Здесь $\mathcal H$~--- допустимые конфигурации классификатора, модели источника и политики, $L_{\mathrm{sys}}$~--- задержка отклика системы, а $L_{\max}$~--- её допустимая граница. Значение $\lambda$ и граница задержки фиксируются до сравнения конфигураций. Внешние критерии~--- точность P300 и конечного символа, символы в минуту, частота исправлений и задержка отклика; они сообщаются отдельно, поскольку одно значение $Q_{\mathcal D_{\pi}}$ не описывает все практические свойства спеллера.

Для оценки без утечки контексты, фразы и ЭЭГ одного участника не должны попадать одновременно в обучающую и контрольную части соответствующего протокола. Обучение и подбор параметров проводятся без доступа к контрольным меткам. Офлайн-записи позволяют сопоставлять модели источника при совместимых стимулах; влияние динамического предъявления на сигнал и реальную скорость требует онлайн-проверки.

% Разделы метода, экспериментов, обсуждения и заключения сохраняются как
% места для подтверждённых материалов после выполнения исследования.


%%%%%%%%%%%%%%%%%%%%%%%%%%%%%%%%%%%%%%%%%%%%%%%%%%%%%%%%%%%%

\bibliographystyle{unsrtnat}
\bibliography{references}

%%%%%%%%%%%%%%%%%%%%%%%%%%%%%%%%%%%%%%%%%%%%%%%%%%%%%%%%%%%%

\end{document}
